\documentclass[11pt,a4paper]{article}

% --- Encoding ---
\usepackage[utf8]{inputenc}
\usepackage[T1]{fontenc}
\usepackage[spanish,es-noshorthands]{babel}
\usepackage{lmodern}

% --- Layout ---
\usepackage[margin=2.5cm]{geometry}
\usepackage{parskip}
\setlength{\parindent}{0pt}
\setlength{\parskip}{0.5em plus 0.2em minus 0.1em}

% --- Tables and graphics ---
\usepackage{booktabs}
\usepackage{array}
\usepackage{enumitem}
\usepackage{xcolor}
\usepackage{graphicx}
\usepackage{tikz}
\usetikzlibrary{positioning,shapes.geometric,arrows.meta,fit,backgrounds,calc,decorations.pathreplacing}
\usepackage{caption}
\usepackage{amsmath,amssymb}

% --- Code listings ---
\usepackage{listings}
\definecolor{codegray}{rgb}{0.96,0.96,0.96}
\definecolor{codeborder}{rgb}{0.85,0.85,0.85}
\definecolor{codecomment}{rgb}{0.30,0.55,0.30}
\definecolor{codekeyword}{rgb}{0.10,0.30,0.70}
\definecolor{codestring}{rgb}{0.65,0.15,0.10}

\lstdefinestyle{cstyle}{
    language=C,
    backgroundcolor=\color{codegray},
    basicstyle=\ttfamily\footnotesize,
    breaklines=true,
    frame=single,
    rulecolor=\color{codeborder},
    showstringspaces=false,
    keywordstyle=\color{codekeyword}\bfseries,
    commentstyle=\color{codecomment}\itshape,
    stringstyle=\color{codestring},
    aboveskip=1em,
    belowskip=1em,
}

\lstdefinestyle{shellstyle}{
    language=bash,
    backgroundcolor=\color{codegray},
    basicstyle=\ttfamily\footnotesize,
    breaklines=true,
    frame=single,
    rulecolor=\color{codeborder},
    showstringspaces=false,
    keywordstyle=\color{codekeyword}\bfseries,
    commentstyle=\color{codecomment}\itshape,
    stringstyle=\color{codestring},
    aboveskip=1em,
    belowskip=1em,
    columns=fullflexible,
}

\lstdefinestyle{textstyle}{
    backgroundcolor=\color{codegray},
    basicstyle=\ttfamily\footnotesize,
    breaklines=true,
    frame=single,
    rulecolor=\color{codeborder},
    showstringspaces=false,
    aboveskip=1em,
    belowskip=1em,
    columns=fullflexible,
}

\lstset{style=shellstyle}

% --- Definition / Idea / Important boxes ---
\usepackage[most]{tcolorbox}
\newtcolorbox{definicion}[1][]{
    colback=blue!5!white,
    colframe=blue!50!black,
    fonttitle=\bfseries,
    title=Definición,
    sharp corners=southwest,
    rounded corners=northwest,
    arc=2pt,
    boxrule=0.6pt,
    left=8pt,right=8pt,top=4pt,bottom=4pt,
    #1
}
\newtcolorbox{idea}[1][]{
    colback=orange!5!white,
    colframe=orange!60!black,
    fonttitle=\bfseries,
    title=Idea intuitiva,
    sharp corners=southwest,
    rounded corners=northwest,
    arc=2pt,
    boxrule=0.6pt,
    left=8pt,right=8pt,top=4pt,bottom=4pt,
    #1
}
\newtcolorbox{importante}[1][]{
    colback=red!5!white,
    colframe=red!50!black,
    fonttitle=\bfseries,
    title=Importante,
    sharp corners=southwest,
    rounded corners=northwest,
    arc=2pt,
    boxrule=0.6pt,
    left=8pt,right=8pt,top=4pt,bottom=4pt,
    #1
}

% --- Hyperlinks ---
\usepackage[colorlinks=true,
            linkcolor=blue!60!black,
            urlcolor=blue!50!black,
            citecolor=blue!60!black]{hyperref}

% --- Color palette ---
\definecolor{sec1}{HTML}{4C72B0}
\definecolor{sec2}{HTML}{DD8452}
\definecolor{sec3}{HTML}{55A868}
\definecolor{sec4}{HTML}{8172B2}
\definecolor{sec5}{HTML}{C44E52}

% --- Title ---
\title{
    \vspace{-2.5em}
    \textbf{Resumen teórico --- Sistemas Operativos} \\[0.3em]
    \large Notas del curso: \emph{Usuarios y seguridad}
}
\author{
    Tomás Rodeghiero \\
    \small Sistemas Operativos --- Universidad Nacional de Río Cuarto
}
\date{2026}

\begin{document}

\maketitle

\begin{abstract}
\noindent
Este documento es un resumen teórico del capítulo \emph{Usuarios y
seguridad} de las \emph{Notas del curso} de Sistemas Operativos
(Marcelo Arroyo, UNRC). Se presentan los mecanismos que un SO debe
ofrecer para soportar múltiples usuarios: la identificación mediante
usuarios y grupos (con sus bases de datos en \texttt{/etc/passwd},
\texttt{/etc/shadow} y \texttt{/etc/group}), el proceso de login y la
autenticación basada en contraseñas con \emph{hash criptográfico} y
\emph{salt} (incluyendo PAM y la autenticación multifactor), el control
de acceso a recursos mediante la \emph{matriz de control de acceso} y
su implementación como ACLs (con el modelo UNIX U/G/O de 9 bits y las
\emph{full ACLs}), los modelos DAC y MAC (con un ejemplo de
Bell-LaPadula), los privilegios de un proceso (real UID, effective
UID, \texttt{setuid}, \texttt{su}, \texttt{sudo}) y las
\emph{capacidades} como mecanismo más fino de control de privilegios.
El objetivo es servir de material de estudio para el final de la
materia.
\end{abstract}

\tableofcontents
\bigskip

% =====================================================================
\section{Introducción}
% =====================================================================

Soportar \textbf{múltiples usuarios} en un sistema operativo no es sólo
una cuestión de comodidad: implica un conjunto de mecanismos que el SO
debe definir e implementar para que los usuarios convivan sin pisarse,
sin invadirse y sin agotar los recursos compartidos. Las Notas listan
siete responsabilidades centrales:

\begin{enumerate}[leftmargin=*]
    \item \textbf{Separación de ambientes de trabajo} de cada usuario
        en el sistema de archivos (cada usuario tiene su propio
        \emph{home}, sus archivos privados, sus configuraciones).
    \item \textbf{Identificar} y \textbf{autenticar} a los usuarios.
    \item Relacionar programas, procesos y usuarios (todo programa que
        se ejecuta corre ``a nombre de'' algún usuario).
    \item Establecer \textbf{límites} en la utilización de recursos
        para que todos los usuarios puedan trabajar.
    \item Ofrecer \textbf{servicios} (comandos y syscalls) de gestión
        de usuarios.
    \item Ofrecer mecanismos de \textbf{control de acceso} a recursos.
    \item Opcionalmente, ofrecer múltiples terminales/consolas (reales
        o virtuales) para que varios usuarios puedan operar en
        paralelo.
\end{enumerate}

\begin{idea}
A grandes rasgos, los temas que siguen tratan: \emph{quién} es un
usuario (sección 2), \emph{cómo} demuestra que es quien dice ser
(sección 3), \emph{a qué} puede acceder (secciones 4 y 5) y con
\emph{qué privilegios} corre un proceso suyo (secciones 6 y 7).
\end{idea}

% =====================================================================
\section{Usuarios, grupos y archivos}
% =====================================================================

Un SO asocia a cada usuario con:

\begin{itemize}[leftmargin=*,nosep]
    \item un \textbf{identificador único} (en UNIX, un entero positivo,
        el \emph{uid}),
    \item un \textbf{rótulo} o \emph{nombre de usuario}, y
    \item un conjunto de \textbf{credenciales de autenticación}.
\end{itemize}

\subsection{Grupos y recursos}

\begin{definicion}
Un \textbf{grupo} representa un conjunto de usuarios. Se usa para
facilitar la gestión del control de acceso: en lugar de listar
permisos usuario por usuario, basta con asignar permisos al grupo y
agregar al usuario al grupo.
\end{definicion}

Un \textbf{recurso} es un archivo de datos o programa, un mecanismo de
IPC (por ejemplo, un pipe), un proceso o un dispositivo de hardware.
El SO debe proteger los recursos respecto a quién puede acceder a
cada uno.

\subsection{El usuario privilegiado}

\begin{importante}
En todo sistema existe \emph{al menos un} \textbf{usuario privilegiado}.
Es quien puede crear nuevos usuarios y asignarles roles o permisos. No
tiene restricciones de acceso a los recursos y actúa como
\textbf{administrador} del sistema.
\begin{itemize}[leftmargin=*,nosep]
    \item En UNIX se llama \textbf{root}.
    \item En MS-Windows se llama \textbf{Administrador}.
\end{itemize}
\end{importante}

\subsection{Gestión de usuarios y grupos}

El sistema ofrece utilidades para crear/eliminar usuarios y grupos.
En UNIX:

\begin{lstlisting}[style=shellstyle]
adduser   alice    # crear usuario
deluser   alice    # eliminar usuario
addgroup  devs     # crear grupo
delgroup  devs     # eliminar grupo
\end{lstlisting}

\subsection{Bases de datos en UNIX}

En los sistemas tipo UNIX, los datos de los usuarios viven en dos
archivos del directorio \verb|/etc|:

\paragraph{\texttt{/etc/passwd}.}
Contiene registros (uno por usuario) de la forma:

\begin{lstlisting}[style=textstyle]
user:x:uid:gid:full name:homedir:shell
\end{lstlisting}

\begin{itemize}[leftmargin=*,nosep]
    \item \texttt{user}: nombre de usuario.
    \item \texttt{x}: si está presente, significa que la contraseña
        está en \verb|/etc/shadow| (separación deliberada por
        seguridad).
    \item \texttt{uid}, \texttt{gid}: identificadores numéricos de
        usuario y grupo principal.
    \item \texttt{full name}: nombre completo (campo descriptivo).
    \item \texttt{homedir}: directorio \emph{home}.
    \item \texttt{shell}: shell que se le abre al hacer login (por
        ejemplo \verb|/bin/bash|).
\end{itemize}

\paragraph{Usuarios sin login.}
Una configuración típica incluye usuarios que \emph{no} tienen acceso
interactivo: usan \verb|/usr/bin/nologin| como shell. Ejemplos:
\texttt{bin}, \texttt{sys}, \texttt{mail}. Permiten asociar un
\emph{dueño}, un grupo y permisos a programas o servicios del sistema
sin que nadie pueda iniciar sesión con esa cuenta.

\paragraph{Identificadores especiales.}
\begin{itemize}[leftmargin=*,nosep]
    \item \texttt{root} tiene \texttt{uid = 0}.
    \item Al primer usuario \emph{real} del sistema se le asigna
        comúnmente \texttt{uid = 1000}.
\end{itemize}

\paragraph{\texttt{/etc/group}.}
Cada usuario tiene un grupo \emph{principal} y puede ser miembro de
otros. La base de grupos está en \verb|/etc/group|.

\begin{idea}
La idea de poner la contraseña en \verb|/etc/shadow| (legible sólo por
root) y dejar \verb|/etc/passwd| legible por todo el mundo es lo que
permite a comandos comunes (\texttt{ls -l}, \texttt{ps}, \ldots)
traducir uids a nombres sin exponer los hashes de las contraseñas.
\end{idea}

% =====================================================================
\section{Login y autenticación}
% =====================================================================

El sistema ejecuta en cada terminal el programa \texttt{login}, que
generalmente pide el nombre de usuario y las credenciales de
autenticación. El método más común es la presentación de una
\textbf{contraseña} (\emph{password}).

\begin{itemize}[leftmargin=*,nosep]
    \item Cada usuario puede cambiar su contraseña con \texttt{passwd}
        (UNIX).
    \item El administrador asigna la contraseña inicial al crear el
        usuario.
\end{itemize}

\subsection{Formato de \texttt{/etc/shadow}}

Las contraseñas se almacenan en \verb|/etc/shadow|. Cada línea tiene
el formato:

\begin{lstlisting}[style=textstyle]
user:hashed-password:last-changed:min-days:max-days:warn-days:ina...
\end{lstlisting}

\begin{itemize}[leftmargin=*]
    \item \texttt{user}: nombre del usuario.
    \item \texttt{hashed-password}: \textbf{hash de la contraseña},
        comúnmente prefijado por \verb|$n$|, donde $n$ indica el código
        del algoritmo de hash utilizado.
    \item Los demás campos contienen fechas o días desde el último
        cambio, período en el que el sistema debería pedir el cambio
        de contraseña, días de inactividad y la fecha de expiración de
        la cuenta.
\end{itemize}

\subsection{Hash criptográfico}

Para no almacenar el secreto en forma plana, se guarda
$\textit{hash}(\textit{passwd})$. La función hash es un \textbf{hash
criptográfico} (MD5, SHA, \ldots).

\begin{definicion}
Una función \textbf{hash criptográfica} es una función
$\textit{hash} : \{0,1\}^* \to \{0,1\}^n$, donde $n$ es el número de
bits de la salida. Toma una entrada de longitud arbitraria y produce
una salida de longitud fija. Por ejemplo, MD5 produce 128\,bits.
\end{definicion}

\paragraph{Propiedad clave.}
Una de las propiedades deseables es la \textbf{resistencia a
preimágenes}: dado el par $(x,\;\textit{hash}(x))$, debe ser
\emph{muy} difícil encontrar un $x'$ tal que
$\textit{hash}(x') = \textit{hash}(x)$. Esto requiere gran poder de
cómputo por ser un problema NP-Hard.

\begin{importante}
Actualmente esa propiedad está \emph{quebrada} para MD5: es
relativamente fácil romperla con el cómputo moderno, por lo que MD5 no
debe usarse para autenticación de contraseñas.
\end{importante}

\subsection{Verificación en el login}

El programa \texttt{login} autentica al usuario verificando que:
\[
\textit{hash}(\textit{provided\_password}) \;=\; \textit{saved\_password}.
\]
Es decir, nunca compara contraseñas en texto plano; sólo hashes.

\subsection{Salt y ataques por diccionario}

Comúnmente, \texttt{passwd} agrega al password provisto un prefijo
llamado \textbf{salt}: un valor generado aleatoriamente (en UNIX,
típicamente dos caracteres) que se almacena como prefijo del hash.

\paragraph{Por qué hace falta.}
El salt previene los \textbf{ataques por diccionario}.

\begin{idea}
Un ataque por diccionario aplica cuando un atacante \emph{roba} el
archivo de contraseñas. El atacante dispone de un \emph{diccionario}
generado previamente: un conjunto de pares
$(\textit{plaintext},\;\textit{hash}(\textit{plaintext}))$ para
millones de palabras y combinaciones comunes. Si una contraseña (hash)
robada coincide con el segundo elemento de algún par del diccionario,
el atacante ya conoce el texto plano correspondiente: la inferencia se
reduce a una búsqueda en el diccionario.

El \textbf{salt} hace que para cada usuario se calcule el hash con un
prefijo distinto. Para que el ataque sirva, el diccionario debería
contener \emph{todas} las combinaciones (\emph{plaintext + salt});
generarlo se vuelve incomparablemente más costoso.
\end{idea}

\subsection{Otros mecanismos de autenticación}

El sistema de login puede aceptar otras credenciales:

\begin{itemize}[leftmargin=*]
    \item \textbf{Biométricos} en sistemas físicos: huellas dactilares,
        reconocimiento de rostro, etc.
    \item \textbf{Autenticación multifactor.} Por ejemplo, doble
        factor: el usuario provee, además, un código enviado por un
        servicio de mensajería (mail, chat, \ldots) o generado por un
        generador de códigos de autenticación.
\end{itemize}

\subsection{PAM (Pluggable Authentication Modules)}

GNU/Linux ofrece un sistema de autenticación muy flexible llamado
\textbf{Pluggable Authentication Modules (PAM)}. PAM permite
\emph{configurar}, \emph{usar} y \emph{desarrollar} mecanismos de
autenticación que pueden combinarse (por ejemplo: contraseña +
huella + token).

% =====================================================================
\section{Control de acceso}
% =====================================================================

El sistema debe asociar \textbf{permisos de acceso} (lectura,
escritura o ejecución) entre los usuarios y los recursos (archivos de
datos y programas).

\subsection{Matriz de control de acceso}

Conceptualmente, esto se modela mediante una \textbf{matriz de control
de acceso}: las filas son los recursos, las columnas son los usuarios,
y cada celda contiene los permisos que el usuario tiene sobre el
recurso. Por ejemplo:

\begin{center}
\small
\begin{tabular}{@{}lccc@{}}
\toprule
\textbf{resources / users} & \textbf{f1} & \textbf{f2} & \textbf{\ldots} \\
\midrule
root  & rwx & rw- & \ldots \\
\ldots & \ldots & \ldots & \ldots \\
alice & r-{}- & -{}-{}- & \ldots \\
\bottomrule
\end{tabular}
\end{center}

\begin{importante}
Implementar esta matriz directamente es un desafío: podría ser
\emph{inmensa} por la cantidad de recursos y usuarios (millones de
archivos por miles de usuarios). Por eso los SO la representan de
forma comprimida, comúnmente por \emph{columnas} (ACLs) o por
\emph{filas} (\emph{capability lists}).
\end{importante}

% =====================================================================
\section{Listas de control de acceso (ACLs)}
% =====================================================================

Una \textbf{ACL} (\emph{Access Control List}) es una representación de
la matriz \textbf{por columnas}: a cada recurso se le asocia el
conjunto $\{(\textit{user},\;\textit{perms}), \ldots\}$ de pares
usuario-permiso.

\subsection{El modelo UNIX: U/G/O}

Para representar este conjunto de forma compacta, en UNIX los usuarios
se clasifican en tres clases:

\begin{enumerate}[leftmargin=*]
    \item \textbf{U (user)}: el dueño (\emph{owner}) del archivo.
    \item \textbf{G (group)}: el grupo (principal) asociado al archivo.
    \item \textbf{O (others)}: el resto de los usuarios del sistema.
\end{enumerate}

Los permisos se asignan al recurso \emph{para cada clase}: permisos
del dueño, permisos del grupo y permisos del resto.

\paragraph{Representación de 9 bits.}
Sólo hacen falta 9 bits: tres para \emph{rwx} del dueño, tres para el
grupo y tres para los otros:
\[
\underbrace{r\,w\,x}_{U} \;\; \underbrace{r\,-\,-}_{G} \;\; \underbrace{r\,-\,-}_{O}
\]

Estos permisos forman parte de los \textbf{metadatos} de cada archivo
(viven en su \emph{inode}).

\subsection{Comandos UNIX}

\begin{lstlisting}[style=shellstyle]
chown alice file        # cambiar dueno
chgrp devs  file        # cambiar grupo
chmod 644   file        # cambiar permisos (rw-r--r--)
\end{lstlisting}

\subsection{Limitación y full-ACLs}

\begin{idea}
La agrupación en sólo tres clases a veces \emph{no} provee la
flexibilidad necesaria: por ejemplo, dar lectura a Alice y Bob pero no
a otros miembros del grupo no es expresable con 9 bits.
\end{idea}

Los SO modernos ofrecen \textbf{full-ACLs}, que permiten asociar el
conjunto de pares (usuario, permisos) de manera \emph{explícita}. Se
implementan comúnmente como archivos especiales (ocultos) asociados a
cada archivo. En GNU/Linux:

\begin{lstlisting}[style=shellstyle]
getfacl file            # listar ACL extendida
setfacl -m u:alice:rw file   # otorgar rw a alice
\end{lstlisting}

\subsection{DAC vs.\ MAC}

\begin{definicion}
Un sistema en el que los usuarios pueden modificar los permisos de sus
recursos (de los cuales son dueños) se denomina \textbf{discretional
access control (DAC)}. UNIX por defecto es DAC: el dueño decide quién
puede leer/escribir/ejecutar su archivo.
\end{definicion}

Existen modelos en los que los permisos son asignados por una
\textbf{autoridad central de seguridad}, no por el dueño. Se conocen
como \textbf{mandatory access control (MAC)} y son comunes en
organizaciones militares.

\paragraph{Ejemplo: Bell-LaPadula.}
Asigna a los usuarios un \emph{nivel de seguridad} y a los recursos
un \emph{rótulo}. Existe un orden parcial (retículo) de niveles:
\[
\textit{TOP\;SECRET} > \textit{SECRET} > \textit{CONFIDENTIAL} > \textit{PUBLIC}.
\]
Un general con nivel \emph{SECRET} podrá acceder a recursos con
rótulo \emph{SECRET}, \emph{CONFIDENTIAL} o \emph{PUBLIC}, pero
\emph{no} podrá acceder a recursos \emph{TOP SECRET}.

% =====================================================================
\section{Privilegios de un proceso}
% =====================================================================

Un proceso usa los permisos de su \textbf{usuario real} (\emph{real
UID}): el usuario que lanzó su ejecución. Esto le permite acceder a
los recursos permitidos a ese usuario.

\subsection{Effective UID y \texttt{setuid}}

\begin{definicion}
El \textbf{effective UID} (\emph{euid}) de un proceso corresponde al
usuario cuyos permisos efectivamente usa el kernel para validar
accesos. En general coincide con el real UID. Pero, si un programa
tiene activado el permiso especial \textbf{setuid}, al ejecutarse el
\emph{euid} del proceso pasa a ser el del \emph{dueño del programa},
no el del usuario que lo invocó.
\end{definicion}

\paragraph{Caso típico: \texttt{passwd}.}
El programa \verb|/usr/bin/passwd| tiene dueño \texttt{root} y requiere
acceso a \verb|/etc/shadow|, el cual sólo es legible/escribible por
root (para evitar que cualquier usuario pueda copiarlo y atacarlo).
Pero cualquier usuario debe poder cambiar su propia contraseña.

\begin{itemize}[leftmargin=*,nosep]
    \item \texttt{passwd} tiene activado el bit \textbf{setuid}.
    \item Cuando Alice lo ejecuta, el proceso adquiere los permisos de
        \texttt{root}.
    \item De ese modo, el proceso puede actualizar \verb|/etc/shadow|
        a pesar de haber sido lanzado por Alice.
\end{itemize}

\begin{importante}
Esto es un mecanismo de \textbf{escalado de privilegios}: pasar
temporalmente a un nivel de privilegios mayor para realizar una
operación concreta. Es poderoso pero peligroso: cualquier bug en un
programa setuid-root puede ser explotado para obtener root completo.
\end{importante}

\subsection{Otros mecanismos: \texttt{su} y \texttt{sudo}}

Comandos que permiten cambios de privilegio bajo control:

\begin{itemize}[leftmargin=*]
    \item \texttt{su}: permite crear una nueva sesión (\emph{shell})
        con otro usuario (por omisión, \texttt{root}). \texttt{su}
        solicita las credenciales de autenticación del \emph{otro}
        usuario; éste debe tener permitido el login.
    \item \texttt{sudo}: permite ejecutar un comando como otro
        usuario. Sólo requiere reingresar las credenciales del
        \emph{usuario actual}. Su configuración vive en
        \verb|/etc/sudoers| y especifica qué usuarios pueden ejecutar
        qué comandos bajo los permisos de qué otros usuarios.
\end{itemize}

\begin{idea}
La diferencia clave es de quién se piden las credenciales: \texttt{su}
exige conocer la contraseña del usuario destino (típicamente root),
mientras que \texttt{sudo} confía en el usuario actual previa
configuración en \verb|/etc/sudoers|. Por eso \texttt{sudo} se
prefiere en sistemas multi-administrador: no hay que compartir la
contraseña de root.
\end{idea}

En MS-Windows el comando equivalente a \texttt{su} es \texttt{runas}.

% =====================================================================
\section{Capacidades}
% =====================================================================

Que un proceso adquiera privilegios de \texttt{root} (\emph{todos} los
permisos) puede generar problemas de seguridad: requiere que el
programa no contenga errores y que acceda \emph{sólo} a los recursos
estrictamente necesarios.

\subsection{Qué son las capacidades}

\begin{definicion}
Una \textbf{capacidad} (\emph{capability}) define qué operaciones se
pueden aplicar sobre un recurso. Comúnmente se usan para
\emph{limitar} el conjunto de operaciones posibles, en lugar de
otorgar privilegios totales.
\end{definicion}

Por ejemplo, si se abre un archivo en modo de sólo lectura, el
descriptor contiene sólo la capacidad de lectura, inhibiendo cualquier
operación de escritura o truncado.

\subsection{Capacidades como reemplazo de setuid-root}

Muchos SO modernos usan capacidades en programas privilegiados para
limitar el conjunto de operaciones o recursos a los que pueden
acceder.

\paragraph{Ejemplo: \texttt{tcpdump}.}
El programa \texttt{tcpdump} permite \emph{capturar} paquetes que
pasan por una interfaz de red. Para esto requeriría privilegios de
\texttt{root}. La solución moderna es aplicar una \textbf{capacidad}
que sólo permita acceder a \emph{raw sockets}: en GNU/Linux esa
capacidad es \verb|CAP_NET_RAW|.

\begin{idea}
En lugar de darle a \texttt{tcpdump} ``poder absoluto'' como
setuid-root, se le otorga \emph{exactamente} la capacidad que
necesita. Si el binario tuviera un bug explotable, el atacante
ganaría \verb|CAP_NET_RAW|, no root completo: un daño mucho menor.
\end{idea}

Para más información sobre capacidades en GNU/Linux:

\begin{lstlisting}[style=shellstyle]
man capabilities
\end{lstlisting}

% =====================================================================
\section{Cierre}
% =====================================================================

El subsistema de usuarios y seguridad implementa, en conjunto, la idea
de que un sistema multi-usuario debe \emph{identificar} a sus usuarios,
\emph{autenticar} su identidad y \emph{autorizar} cada operación que
intentan realizar. Las ideas principales:

\begin{itemize}[leftmargin=*]
    \item Cada usuario tiene un \textbf{uid} y pertenece a uno o más
        \textbf{grupos}. La separación de \texttt{/etc/passwd},
        \texttt{/etc/shadow} y \texttt{/etc/group} es la concreción de
        este modelo en UNIX. Root (\texttt{uid = 0}) es el único
        usuario sin restricciones.
    \item La \textbf{autenticación} compara hashes, no texto plano.
        Una función hash criptográfica garantiza que conocer
        $\textit{hash}(x)$ no permite recuperar $x$ con facilidad. El
        \textbf{salt} es la defensa específica contra ataques por
        diccionario, multiplicando exponencialmente el costo de
        precomputar tablas. PAM permite encadenar varios mecanismos
        (contraseña, biométricos, multifactor).
    \item El \textbf{control de acceso} se piensa con una matriz
        recurso $\times$ usuario, pero se implementa como ACLs por
        columna. UNIX comprime la ACL en 9 bits con las clases
        \emph{U/G/O}, complementadas con \emph{full ACLs} cuando hace
        falta más granularidad.
    \item Hay dos filosofías: \textbf{DAC} (el dueño decide; UNIX por
        defecto) y \textbf{MAC} (una autoridad central decide; típico
        de entornos militares, con modelos como Bell-LaPadula).
    \item Cada proceso tiene un \emph{real UID} y un
        \textbf{effective UID}. El bit \textbf{setuid} permite
        ejecutar un programa con los privilegios de su dueño,
        habilitando casos legítimos como \texttt{passwd} pero abriendo
        riesgos de escalado. \texttt{su} y \texttt{sudo} ofrecen
        cambios de privilegio bajo control.
    \item Las \textbf{capacidades} reemplazan ``setuid-root'' por
        privilegios \emph{finos y específicos}: en lugar de dar todo
        el poder, dar exactamente lo necesario (por ejemplo,
        \texttt{CAP\_NET\_RAW} a \texttt{tcpdump}).
\end{itemize}

\begin{idea}
El principio rector que atraviesa todo el capítulo es el de
\textbf{mínimo privilegio}: cada usuario, proceso y programa debe
tener exactamente los permisos que necesita, ni más ni menos. Las
ACLs, \texttt{sudo}, las capacidades y los modelos MAC son técnicas
distintas para acercarse a ese ideal en escenarios reales.
\end{idea}

\end{document}
